\section{Introduction}
\label{sec:introduction}



The pursuit of creating intelligent and autonomous agents that can seamlessly assist humans and operate in real-world settings has been a foundational goal in artificial intelligence~\citep{wooldridge1995intelligent,Minsky88,bubeck2023sparks}. 
The recent advance of Large Language Models (LLMs)~\citep{DBLP:journals/corr/abs-2303-08774,DBLP:journals/corr/abs-2305-10403,DBLP:journals/corr/abs-2307-09288} has created newfound avenues in this domain. These LLMs, especially GPT-4 \citep{DBLP:journals/corr/abs-2303-08774}, are particularly adept in comprehending human intent and executing commands. They have demonstrated remarkable proficiency in domains such as language understanding, vision~\citep{OpenAI_ChatGPT}, and coding~\citep{bubeck2023sparks}. 
By harnessing the power of LLMs, autonomous agents can make more nuanced decisions and perform actions with an unprecedented degree of autonomy \citep{DBLP:journals/corr/abs-2307-13854}. Agents like AutoGPT \citep{richards2023auto}, BabyAGI \citep{nakajima2023babyagi}, and AgentGPT \citep{agentgpt}, are inspiring examples. Furthermore, recent research has endowed autonomous agents with more human-analogous cognitive mechanisms, spanning from reflection \citep{yao2023react,shinn2023reflexion}, task decomposition \citep{wei2023chainofthought,yao2023tree}, and tool utilization \citep{schick2023toolformer,qin2023tool,qin2023toolllm,qian2023creator}. These advancements edge us closer to realizing the concept of artificial general intelligence (AGI) \citep{goertzel2007artificial,clune2020aigas} that can generalize across a broader range of tasks.




However, complex real-world tasks often require cooperation among individuals to achieve better effectiveness. Throughout history, numerous studies have delved into methods for enhancing collaboration among humans to improve work efficiency and effectiveness~\citep{collectiveintelligence,fehr2000cooperation}. More recently, with the evolution of autonomous agents towards AGI, extensive research conceptualizes the assemblies of agents as a society or group \citep{DBLP:journals/corr/abs-2303-17760}, and focuses on exploring the potential of their cooperation. For example, \cite{DBLP:journals/corr/abs-2304-03442} found emergent social behaviors in multi-agent life simulation. \cite{DuMultiagentDebate,wang2023unleashing,zhang2023building,DBLP:journals/corr/abs-2307-07924,chan2023chateval} also underscored the enhanced decision-making of collaborating agents during collaborative problem-solving. 
However, a limitation in these studies is their narrow focus on specific and limited tasks, leaving the generalizability of their findings uncertain. An additional constraint is their static approach to agent collaboration, where agents' roles and capabilities remain rigid, hindering adaptability.



\begin{figure}[!t]
    \centering
    \includegraphics[width=0.95\textwidth]{figs/files/pipeline.pdf}
    \caption{An illustration of the \framework.}
    \label{fig:pipeline}
    \vspace{-1em}
\end{figure}

% 把建立pip的思想說清楚。“把多agent組織成一個有效率的team，並讓他們能在真實環境中更高效的完成任務”。
% 1. 因為多agent比 單agent好已經有人研究過
% 2. 對話框架這個在 CahtDev, ChatEval, Heng Ji's 文章說過
% 3. 與環境交互 單 Agent 的文章也提過
% -- 所以應該定為在於說，提出pipeline，這pipeline能把多Agents，建立一個有效率的Team，這個team的人員可以 “依據當前的環境動態變化(替換)，去調整組織，以更好地做出決策”
% 因此我們提出AgentPip - 主要模擬一個Group在accomplish Goal時的過程，這過程需要，1.諮詢新的人.. 2.

To address this problem, we introduce \framework. This general multi-agent framework simulates the problem-solving procedures of human groups, and allows for dynamic adjustment of group members based on current progress.
Specifically, \framework splits the problem-solving process into four pivotal stages as shown in \cref{fig:pipeline}:
(1) \textit{Expert Recruitment}: Determine and adjust the agent group's composition based on the ongoing problem-solving progression.
(2) \textit{Collaborative Decision-Making}: Engage the selected agents in joint discussions to devise problem-solving strategies. 
(3) \textit{Action Execution}: Agents interact with their environment to implement the devised actions.
(4) \textit{Evaluation} - Assess the differences between the current state and desired outcomes. If the current state is unsatisfactory, feedback is given to the next iteration for further refinement.


We conduct extensive experiments and case studies in diverse aspects including text understanding, reasoning, coding, tool utilization and embodied AI to show the effectiveness of \framework. Additionally, we highlight the social behaviors that emerge from the multi-agent collaboration, and discuss their advantages and potential risks. In summary, our contributions are:

\begin{itemize}[noitemsep,topsep=0pt,parsep=0pt,partopsep=0pt,leftmargin=2em]
    \item Inspired by the collaborative process of a human team, we propose \framework as \textit{an effective framework for promoting collaboration among multiple agents} in problem-solving.
    \item We conduct extensive experiments to show that \framework effectively improve the agents' understanding, reasoning, coding, tool utilizing capabilities and their potential in embodied AI.
    \item In the multi-agent collaboration, especially within tool utilization and Minecraft game playing, agents manifest certain emergent behaviors. For example, (1) \textit{volunteer behaviors}, characterized by agents offering assistance to peers, thus improving team efficiency; (2) \textit{conformity behaviors}, where agents adjust their deviated behaviors to align with the common goal under the critics from others; (3) \textit{destructive behaviors}, occasionally leading to undesired and detrimental outcomes. 
\end{itemize}


